% !TeX program = pdflatex
\documentclass[12pt,a4paper]{article}

% Encoding & sprog
\usepackage[utf8]{inputenc}
\usepackage[T1]{fontenc}
\usepackage[danish]{babel}

% Sideopsætning
\usepackage{geometry}


% Grafik
\usepackage{graphicx}

% Header/footer
\usepackage{graphicx} % i preamble
\usepackage{lastpage}
\usepackage{fancyhdr}
\usepackage{enumerate}
\usepackage{mathtools}
\usepackage{amsthm}
\usepackage{pgf}
\usepackage{caption}
\usepackage{subcaption}
\usepackage{blkarray}
\usepackage{qtree}
\usepackage{amsfonts}
\usepackage[shortlabels]{enumitem}
\usepackage[utf8]{inputenc}
\usepackage{stmaryrd}
\usepackage{physics}
\usepackage{amsmath}
\usepackage{tikz}
\usepackage{amssymb}
\usepackage{listings}
\usepackage{xcolor}
\usepackage{minted}
\usepackage{xcolor} % to access the named colour LightGray
\definecolor{LightGray}{gray}{0.9}

\setminted[python]{ %minted design :3
  frame=lines,
  framesep=2mm,
  baselinestretch=0.5,
  bgcolor=LightGray,
  fontsize=\footnotesize,
  linenos
}



\renewcommand{\thesubsection}{\thesection.\alph{subsection}}
\geometry{top=3.5cm, bottom=3cm, left=3cm, right=3cm, headheight=2.5cm}
\pagestyle{fancy}
\fancyhf{}
\lhead{\includegraphics[width=10cm]{Logo/aulogo.png}\\}
\chead{\raisebox{6mm}{Digital Signal Behandling - HandIn 2}}
\rhead{\raisebox{6mm}{13/10/2026}}
\rfoot{Side \thepage\ af \pageref{LastPage}}
\renewcommand{\headrulewidth}{0pt}

% Hyperref (indlæses altid sidst)
\usepackage[hidelinks]{hyperref}

\begin{document}

% ----------- FORSIDE -----------
\pagenumbering{gobble}
\pagestyle{empty}
\begin{titlepage}
    \thispagestyle{empty}
    \centering

    {\Large \textsc{Aarhus University}} \\
    \vspace{2.5cm}

    {\Huge \textbf{Frekvensanalysesystem baseret
på Diskret Fourier Transformation}} \\
    \vspace{0.5cm}

    {\Large 13/10/2026} \\
    \vspace{2cm}

    {\large \textit{Group X}} \\
    \vspace{0.2cm}

    {\large Yusuf \textsc{Cakir} - AU-ID.: AU802700} \\
    {\large Firstname \textsc{Lastname} - AU-ID.: XXXXXXXX}\\
    {\large Asger Grundtdal \textsc{Stidsen} - AU-ID.: AU798909}
    \\

    \vspace{2cm}

    {\large \textit{Instructors}} \\
    {\large Lars \textsc{G. Johansen}} \\


    \vfill

    \includegraphics[width=0.45\textwidth]{Logo/aulogoblue.png}

\end{titlepage}
\setcounter{page}{1}
%AUTO indholdfortegnelse
\tableofcontents

% ----------- INDHOLD -----------
\newpage
\pagenumbering{arabic}
\setcounter{page}{1}
\pagestyle{fancy}



\section*{Exercise 1} %1
De første 10 sekunder af begge støjsignaler analyseres.
For et samplet signal med samplingsfrekvens \(f_s\) er samplingstiden:
\[
f_s = \frac{1}{T_s}
\]
Tidsaksen dannes som:
\[
t[n]=\frac{n}{f_s}
\]
hvor:
\[
n=0,1,\ldots,N-1
\]
Begge lydfiler har:
\[
f_s=48.000\text{ Hz}
\]
For et signal på 10 sekunder bliver antallet af samples:
\[
N=T\cdot f_s
\]
\[
N=10\cdot48.000
\]
\[
N=480.000
\]


\section*{Exercise 2} %2

\section*{Exercise 3} %3

\section*{Exercise 4} %4

\code[py]{}{code:e4-}
\begin{lstlisting}


\end{lstlisting}


\section*{Exercise 5} %5

\section*{Exercise 6} %6

\section*{Exercise 7} %7
 
\section*{Exercise 4.13}

\begin{figure}[H]
    \centering
    \includegraphics[width=1\linewidth]{Ekstraopgave/13.png}
    \caption{Enter Caption}
    \label{fig:placeholder}
\end{figure}

\begin{minted}{python}
fs = 2000
required_df = 0.5

N_min = fs / required_df
N_fft = 2**np.ceil(np.log2(N_min))   # FFT kræver N som potens af 2
df_actual = fs / N_fft

print("Minimum N:", N_min)
print("N (potens af 2):", int(N_fft))
print("Reel frekvensopløsning:", df_actual, "Hz")
\end{minted}

\begin{table}[h]
\centering
\renewcommand{\arraystretch}{1.3}
\begin{tabular}{l r}
\hline
\textbf{Størrelse} & \textbf{Værdi} \\
\hline
Minimum $N$ & 4000.0 \\
$N$ (potens af 2) & 4096 \\
Reel frekvensopløsning & 0.48828125 Hz \\
\hline
\end{tabular}
\caption{Bestemmelse af antal datapunkter og reel frekvensopløsning}
\end{table}
 
\section*{Exercise 4.17}

\begin{figure}[H]
    \centering
    \includegraphics[width=1\linewidth]{Ekstraopgave/image17.png}
    \caption{Enter Caption}
    \label{fig:placeholder}
\end{figure}


\begin{minted}{python}
N = 10
n = np.arange(N)

# A) Hamming window
w_ham = 0.54 - 0.46*np.cos(2*np.pi*n/(N-1))

# B) Hanning window
w_han = 0.5 - 0.5*np.cos(2*np.pi*n/(N-1))

print("A) Hamming:", np.round(w_ham, 4))
print("B) Hanning:", np.round(w_han, 4))
\end{minted}


\begin{table}[h]
\centering
\resizebox{\textwidth}{!}{%
\begin{tabular}{c|cccccccccc}
\hline
$n$ & 0 & 1 & 2 & 3 & 4 & 5 & 6 & 7 & 8 & 9 \\
\hline
A) Hamming & 0.08 & 0.1876 & 0.4601 & 0.77 & 0.9723 & 0.9723 & 0.77 & 0.4601 & 0.1876 & 0.08 \\
B) Hanning & 0 & 0.117 & 0.4132 & 0.75 & 0.9698 & 0.9698 & 0.75 & 0.4132 & 0.117 & 0 \\
\hline
\end{tabular}%
}
\caption{Hamming- og Hanning-vinduesfunktioner for $N=10$}
\end{table}

 
\section*{Exercise 4.18}
 
 \begin{figure}[H]
     \centering
     \includegraphics[width=1\linewidth]{Ekstraopgave/image2.png}
     \caption{Enter Caption}
     \label{fig:placeholder}
 \end{figure}

\begin{minted}{python}

x = np.array([0, 0.2, 0, -0.2, 0, 0.2])
N = 6
n = np.arange(N)

# A) Triangular window
w_tri = 1 - np.abs((n - (N-1)/2) / ((N-1)/2))
xw_tri = x * w_tri

# B) Hamming window
w_ham = 0.54 - 0.46*np.cos(2*np.pi*n/(N-1))
xw_ham = x * w_ham

# C) Hanning window
w_han = 0.5 - 0.5*np.cos(2*np.pi*n/(N-1))
xw_han = x * w_han

print("A) Triangular:")
print("  w(n)  =", np.round(w_tri, 4))
print("  xw(n) =", np.round(xw_tri, 4))

print("\nB) Hamming:")
print("  w(n)  =", np.round(w_ham, 4))
print("  xw(n) =", np.round(xw_ham, 4))

print("\nC) Hanning:")
print("  w(n)  =", np.round(w_han, 4))
print("  xw(n) =", np.round(xw_han, 4))

\end{minted}

\begin{table}[h]
\centering
\renewcommand{\arraystretch}{1.3}
\begin{tabular}{c|cccccc}
\hline
$n$ & 0 & 1 & 2 & 3 & 4 & 5 \\
\hline
\multicolumn{7}{l}{\textbf{A) Triangular window}} \\
\hline
$w(n)$  & 0 & 0.4 & 0.8 & 0.8 & 0.4 & 0 \\
$x_w(n)$ & 0 & 0.08 & 0 & $-0.16$ & 0 & 0 \\
\hline
\multicolumn{7}{l}{\textbf{B) Hamming window}} \\
\hline
$w(n)$  & 0.08 & 0.3979 & 0.9121 & 0.9121 & 0.3979 & 0.08 \\
$x_w(n)$ & 0 & 0.0796 & 0 & $-0.1824$ & 0 & 0.016 \\
\hline
\multicolumn{7}{l}{\textbf{C) Hanning window}} \\
\hline
$w(n)$  & 0 & 0.3455 & 0.9045 & 0.9045 & 0.3455 & 0 \\
$x_w(n)$ & 0 & 0.0691 & 0 & $-0.1809$ & 0 & 0 \\
\hline
\end{tabular}
\caption{Vindues- og vinduede sekvenser for $N=6$}
\end{table}
 
 \section*{Exercise 4.20}

\begin{figure}[H]
     \centering
     \includegraphics[width=1\linewidth]{Ekstraopgave/image.png}
     \caption{Enter Caption}
     \label{fig:placeholder}
 \end{figure}

 \begin{minted}{python}

fs = 8000   # samplingsfrekvens
f0 = 2000   # signalfrekvens

# A
N = 100
df = fs / N
ratio = f0 / df
leakage = "Nej" if ratio == int(ratio) else "Ja"
print(f"A) N = {N}: Δf = {df:.4f} Hz, f0/Δf = {ratio:.4f}, Lækage: {leakage}")

# B
N = 73
df = fs / N
ratio = f0 / df
leakage = "Nej" if ratio == int(ratio) else "Ja"
print(f"B) N = {N}: Δf = {df:.4f} Hz, f0/Δf = {ratio:.4f}, Lækage: {leakage}")
\end{minted}


A) N = 100 -> ingen lækage
$\Delta$f = fs/N = 8000/100 = 80 Hz
f0/$\Delta$f = 2000/80 = 25 (heltal)

Da forholdet er et heltal, rammer signalfrekvensen 2000 Hz præcis på et DFT-bin. Signalet indeholder et helt antal perioder inden for de 100 samples, så der opstår ingen diskontinuitet i vinduet, energien samles i ét enkelt bin.

B) N = 73 -> spektral lækage
$\Delta$f = fs/N = 8000/73 = 109.59 Hz
f0/$\Delta$f = 2000/109.59 = 18.25 (ikke heltal)

Da forholdet ikke er et heltal, rammer 2000 Hz mellem to bins. Signalet indeholder ikke et helt antal perioder inden for de 73 samples, hvilket giver en diskontinuitet, når DFT'en behandler segmentet som periodisk gentaget. Denne diskontinuitet spreder energien ud over flere nabo-bins, det er det, der kaldes spektral lækage.

\end{document}